\documentclass[11pt,a4paper]{article}
\usepackage[margin=24mm]{geometry}
\usepackage[T1]{fontenc}
\usepackage{lmodern}
\usepackage{amsmath,amssymb,booktabs,tabularx,array}
\usepackage{xcolor}
\usepackage[unicode,colorlinks=true,linkcolor=blue!45!black,urlcolor=blue!45!black,citecolor=blue!45!black]{hyperref}
\usepackage{enumitem}
\setlist{nosep,leftmargin=*}
\setlength{\emergencystretch}{2em}
\newcommand{\pending}[1]{\textcolor{blue!55!black}{[TO COMPLETE: #1]}}
\newcommand{\unmeasured}{\textit{not measured}}
\newcommand{\projectstatus}{\begin{quote}\small\textbf{Research planning document, 25 November 2026.}
This document specifies a proposed study. Empty results are intentional. No experiment,
human annotation, acceptance, or publication is implied by its prose.
The eventual title, abstract, claims, and discussion must be rewritten from verified evidence.\end{quote}}
\newcolumntype{Y}{>{\raggedright\arraybackslash}X}

\title{Tokens or State Changes? A Controlled Study of Delayed Supervision for Compact Language Models}
\author{}
\date{Research design: 25 November 2026}
\begin{document}
\maketitle
\projectstatus
\begin{abstract}
This planned study separates token delay from semantic-update interference in
training compact language models. Four supervision schedules share event-language
data and matched answer-label budgets: immediate probes, delayed probes with few
intervening state changes, matched token delays with many state changes, and a
fixed mixture. Evaluation asks whether any benefit transfers from short probe
answers to coherent multi-sentence continuation. Additional prefix processing,
answer formatting and token exposure are accounted for explicitly.
\pending{Insert actual datasets, model counts, effects, intervals and limits
after validation and registered experiments. No supervision benefit is claimed.}
\end{abstract}

\section{Question and related work}
Delayed supervision already trains later access to simulator-grounded
facts~\cite{delayed}. Its reported comparisons do not isolate delay from QA,
termination supervision and extra computation. State-reveal code traces~\cite{code},
latent-situation supervision~\cite{latent}, permutation mechanisms~\cite{stateicml}
and data-efficiency comparisons across supervision regimes~\cite{bias} further
limit broad novelty claims. We investigate a controlled distinction between
token distance and intervening semantic writes, and whether it matters for
generation rather than only component probes. More seeds or a new backbone
alone would not establish an independent contribution.
SchemaMem also studies selective updates and delayed retrieval in a
schema-indexed recurrent architecture~\cite{schemamem}; our question concerns
the supervision schedule, not a first solution to delayed state access.

\section{Training schedules and label validity}
For each relevant fact, record the last supporting update, token delay
\(D_{\rm tok}\), and number of intervening successful state-changing events
\(D_{\rm write}\). A retention answer must remain valid through the delay;
a revised answer is the latest value, with delay measured from its last relevant
revision. Balance question types and answer classes across conditions.

\begin{table}[ht]\centering\small
\caption{Schedules on a shared event corpus.}
\begin{tabularx}{\linewidth}{lYY}\toprule
Arm & Probe placement & Main control\\\midrule
B0 & Immediately after relevant evidence & QA/answer-format reference\\
B1 & Long token delay, few intervening writes & Distance with low semantic interference\\
B2 & Matched token-delay bands, many other-entity writes & Semantic interference in supervision\\
B3 & Preregistered B1/B2 mixture & Training-schedule diversity\\
\bottomrule\end{tabularx}\end{table}

The event LM corpus remains the same; only probe placement/selection differs.
Question-answer probes run on disposable branches so their answers never enter
the continuing event stream. Use ordinary LM-head answer targets, including a
matched EOS policy, with no gold state table at deployment. The objective is
\[
\mathcal L=\mathcal L_{\rm event}+\lambda\mathcal L_{\rm answer},\qquad
\mathcal L_{\rm answer}= -|\mathcal I|^{-1}\sum_{j\in\mathcal I}
 |a_j|^{-1}\log p_\theta(a_j\mid x_{\le t_j},q_j).
\]
State schema, question inventory and answer targets are controlled, not invented
from a model's prediction. Match answer-target counts and report event-token
exposure. Longer auxiliary prefixes change processed tokens and gradient spans:
they are measured confounds, not free supervision. A measured-compute replication
needs its own predeclared cap. Specify detach/resume behavior; do not silently
truncate the evidence-to-probe gradient and still claim training distant writes.

\section{Common training and evaluation contract}
The output is a three-to-five-sentence autoregressive continuation, not a state
label: \(p_\theta(y\mid x)=\prod_t p_\theta(y_t\mid x,y_{<t})\). Gold event states
are used for scoring and declared training supervision, never as an inference-time
table supplied to the generator. A controlled prefix specifies successful updates,
rule compliance and time/world scope. A point-in-time obligation is evaluated once;
``must act'' entails execution only under the stated compliance assumption.
Possession is distinct from ownership. Multiple coherent narrative choices are valid.

The proposed pilot uses minGRU and a causal Transformer, each with a 512-piece
vocabulary, two layers and three training seeds. The instantiated minGRU at width
604 has \textbf{2,081,384} parameters; the Transformer at width 268 has
\textbf{2,074,320}. These are measured backbone counts, not trained quality results.
The parameter gap is approximately 0.34\%. Count any added heads separately.
The current context of 256 was used for instantiation only; choose a common
context that contains the complete task before training and recount any changes.
Fit the tokenizer on training data only. TinyStories~\cite{tiny} is a candidate
general-English warm-up corpus, subject to version/license and source checks.

Four conditions, two backbones and three seeds give 24 post-training runs. The
planning pilot caps each at two million CE targets (48M total); confirmation at
ten million targets/run (240M) requires prior feasibility. Selection receives a
separate 15\% cap. Six warm-up checkpoints would consume 30M targets at 5M each
in the pilot; report immutable provenance and reuse rather than calling a reused
checkpoint a new run. No token budget guarantees adequate writing. Freeze
data/order/decoding/selection policy within each controlled comparison. Different
auxiliary prefix replays or rollouts require separate processed-token and cost logs.
Parameter matching, target-token matching and measured-compute matching are
different conditions, not interchangeable claims.

For a continuation, let A mean all mandatory propositions are correctly expressed,
F the requested sentence format, and C a contradiction with the applicable prefix
or generated history. Generator success is \(Y=\mathbf1[A\land F\land\neg C]\).
Also report semantic-only success, proposition coverage, changed/preserved facts,
omission/evasion, fluency and held-out perplexity under the same tokenizer.
Silence is not memory success. An unmentioned fact is not automatically false.

Use source-disjoint test stories with a human-authored or appropriately licensed
subset; do not call template-generated naturalistic text human-authored. Freeze
the primary contrast, test hashes, sampling and decoding before final selection.
Two readers label a representative sample independently, blinded to system and
condition; retain evidence spans, unclear cases and pre-adjudication disagreement.
No human labels exist yet for this design. Audit automatic judgments against
real labels rather than presuming an oracle parses free text correctly.
Pair source-cluster resampling across conditions for 95\% intervals, retaining
premise variants and repeated outputs together; report every training seed and
seed variability. Sentences are not independent samples. Report naturalistic
transfer separately, not only an aggregate dominated by controlled templates.

Log actual parameters, corpus/tokenizer/model hashes, target and processed tokens,
auxiliary targets, FLOP proxy, peak allocated/reserved VRAM, train/generation time
and selection budget. The observed host is an RTX 5050 Laptop (8,151 MiB).
Throughput for these new configurations, Colab credit burn and total GPU-hours
are unmeasured. No paid API/cloud is assumed. Required implementation checks
include causal masking, story reset, padding, tiny-batch memorization, step/full
eval equivalence and checkpoint resume. The revised objectives/data procedures
are proposals, not implemented by this document.


\section{Primary contrast and transfer}
The primary contrast is B2 minus B1 on continuation success at held-out
delay/write-density conditions. Immediate versus delayed and the mixture are
secondary. Cross token delay with semantic-write load at evaluation, controlling
length and recency. Questions about current states remain component tests;
they do not replace three-to-five-sentence generation.

Report whether improvements survive answer-format controls, naturalistic source
transfer and compute accounting. A separate plot relates probe success to
continuation success without treating correlation as proof of internal storage.
Report optimization failures and confidence intervals, including cases where
probe accuracy improves while writing or semantic coverage does not.

\begin{table}[ht]\centering\small
\caption{Planned evidence table, currently empty.}
\begin{tabularx}{\linewidth}{lYYYY}\toprule
Arm & Continuation / CI & Probe / format & Processed / target tokens & Time / VRAM\\\midrule
B0 & -- & -- & -- & --\\B1 & -- & -- & -- & --\\
B2 & -- & -- & -- & --\\B3 & -- & -- & -- & --\\
\bottomrule\end{tabularx}\end{table}

\section{Limitations and completion criteria}
An effect cannot be called the pure causal effect of time if semantic content,
label validity or processing cost differs. A useful study resolves a named
confound with transfer evidence or a sufficiently precise negative result;
simply adding an auxiliary head is insufficient. Model limitations, supervision
noise, window resets and English-only data restrict conclusions. No new labels,
reliability measurement or trained result exists for this design. Audit probe
branches, source splits, byte/BPE alignment and future-token causality before
final training. Add genuine authorship, data licenses, consent and disclosures,
and rewrite claims from completed evidence.

\begin{thebibliography}{99}
\raggedright
\bibitem{delayed} Jinha Kim and Taksh Kothari. Delayed Supervision for Test-Time Language Models. arXiv preprint, 2026.
\url{https://arxiv.org/abs/2609.32312}.
\bibitem{code} Julien Siems et al. Learning State-Tracking from Code Using Linear RNNs. arXiv version, 2026.
\url{https://arxiv.org/abs/2602.14814}.
\bibitem{bias} Reza Ebrahimi et al. On the Induction Bias in Sequence Models. ICML, 2026.
\url{https://proceedings.mlr.press/v306/ebrahimi26a.html}.
\bibitem{stateicml} Belinda Z. Li, Zifan Carl Guo, and Jacob Andreas. (How) Do Language Models Track State? ICML, 2025.
\url{https://proceedings.mlr.press/v267/li25r.html}.
\bibitem{latent} Language Modeling with Latent Situations. Findings of ACL, 2023.
\url{https://aclanthology.org/2023.findings-acl.795/}.
\bibitem{openmeva} Jian Guan et al. OpenMEVA: A Benchmark for Evaluating Open-ended Story Generation Metrics. ACL, 2021.
\url{https://aclanthology.org/2021.acl-long.500/}.
\bibitem{mingru} Leo Feng et al. Were RNNs All We Needed? arXiv version, 2024.
\url{https://arxiv.org/abs/2410.01201}.
\bibitem{tiny} Ronen Eldan and Yuanzhi Li. TinyStories: How Small Can Language Models Be and Still Speak Coherent English? arXiv preprint, 2023.
\url{https://arxiv.org/abs/2305.07759}.
\bibitem{schemamem} Sungwoo Goo, Hwi-yeol Yun, and Sangkeun Jung.
SchemaMem: Schema-Indexed Recurrent Memory for Delayed State Retrieval.
arXiv preprint, 2026. \url{https://arxiv.org/abs/2609.33436}.
\end{thebibliography}

\end{document}
